\documentclass[a4paper,12pt,twoside]{memoir}

% Castellano
\usepackage[spanish,es-tabla]{babel}
\selectlanguage{spanish}
\usepackage[utf8]{inputenc}
\usepackage[T1]{fontenc}
\usepackage{lmodern} % scalable font
\usepackage{microtype}
\usepackage{placeins}

\RequirePackage{booktabs}
\RequirePackage[table]{xcolor}
\RequirePackage{xtab}
\RequirePackage{multirow}

% Links
\PassOptionsToPackage{hyphens}{url}\usepackage[colorlinks]{hyperref}
\hypersetup{
	allcolors = {red}
}

% Ecuaciones
\usepackage{amsmath}

% Rutas de fichero / paquete
\newcommand{\ruta}[1]{{\sffamily #1}}

% Párrafos
\nonzeroparskip

% Huérfanas y viudas
\widowpenalty100000
\clubpenalty100000

% Evitar solapes en el header
\nouppercaseheads

% Imagenes
\usepackage{graphicx}
\newcommand{\imagen}[2]{
	\begin{figure}[!h]
		\centering
		\includegraphics[width=0.9\textwidth]{#1}
		\caption{#2}\label{fig:#1}
	\end{figure}
	\FloatBarrier
}

\newcommand{\imagenflotante}[2]{
	\begin{figure}%[!h]
		\centering
		\includegraphics[width=0.9\textwidth]{#1}
		\caption{#2}\label{fig:#1}
	\end{figure}
}



% El comando \figura nos permite insertar figuras comodamente, y utilizando
% siempre el mismo formato. Los parametros son:
% 1 -> Porcentaje del ancho de página que ocupará la figura (de 0 a 1)
% 2 --> Fichero de la imagen
% 3 --> Texto a pie de imagen
% 4 --> Etiqueta (label) para referencias
% 5 --> Opciones que queramos pasarle al \includegraphics
% 6 --> Opciones de posicionamiento a pasarle a \begin{figure}
\newcommand{\figuraConPosicion}[6]{%
  \setlength{\anchoFloat}{#1\textwidth}%
  \addtolength{\anchoFloat}{-4\fboxsep}%
  \setlength{\anchoFigura}{\anchoFloat}%
  \begin{figure}[#6]
    \begin{center}%
      \Ovalbox{%
        \begin{minipage}{\anchoFloat}%
          \begin{center}%
            \includegraphics[width=\anchoFigura,#5]{#2}%
            \caption{#3}%
            \label{#4}%
          \end{center}%
        \end{minipage}
      }%
    \end{center}%
  \end{figure}%
}

%
% Comando para incluir imágenes en formato apaisado (sin marco).
\newcommand{\figuraApaisadaSinMarco}[5]{%
  \begin{figure}%
    \begin{center}%
    \includegraphics[angle=90,height=#1\textheight,#5]{#2}%
    \caption{#3}%
    \label{#4}%
    \end{center}%
  \end{figure}%
}
% Para las tablas
\newcommand{\otoprule}{\midrule [\heavyrulewidth]}
%
% Nuevo comando para tablas pequeñas (menos de una página).
\newcommand{\tablaSmall}[5]{%
 \begin{table}
  \begin{center}
   \rowcolors {2}{gray!35}{}
   \begin{tabular}{#2}
    \toprule
    #4
    \otoprule
    #5
    \bottomrule
   \end{tabular}
   \caption{#1}
   \label{tabla:#3}
  \end{center}
 \end{table}
}

%
%Para el float H de tablaSmallSinColores
\usepackage{float}

%
% Nuevo comando para tablas pequeñas (menos de una página).
\newcommand{\tablaSmallSinColores}[5]{%
 \begin{table}[H]
  \begin{center}
   \begin{tabular}{#2}
    \toprule
    #4
    \otoprule
    #5
    \bottomrule
   \end{tabular}
   \caption{#1}
   \label{tabla:#3}
  \end{center}
 \end{table}
}

\newcommand{\tablaApaisadaSmall}[5]{%
\begin{landscape}
  \begin{table}
   \begin{center}
    \rowcolors {2}{gray!35}{}
    \begin{tabular}{#2}
     \toprule
     #4
     \otoprule
     #5
     \bottomrule
    \end{tabular}
    \caption{#1}
    \label{tabla:#3}
   \end{center}
  \end{table}
\end{landscape}
}

%
% Nuevo comando para tablas grandes con cabecera y filas alternas coloreadas en gris.
\newcommand{\tabla}[6]{%
  \begin{center}
    \tablefirsthead{
      \toprule
      #5
      \otoprule
    }
    \tablehead{
      \multicolumn{#3}{l}{\small\sl continúa desde la página anterior}\\
      \toprule
      #5
      \otoprule
    }
    \tabletail{
      \hline
      \multicolumn{#3}{r}{\small\sl continúa en la página siguiente}\\
    }
    \tablelasttail{
      \hline
    }
    \bottomcaption{#1}
    \rowcolors {2}{gray!35}{}
    \begin{xtabular}{#2}
      #6
      \bottomrule
    \end{xtabular}
    \label{tabla:#4}
  \end{center}
}

%
% Nuevo comando para tablas grandes con cabecera.
\newcommand{\tablaSinColores}[6]{%
  \begin{center}
    \tablefirsthead{
      \toprule
      #5
      \otoprule
    }
    \tablehead{
      \multicolumn{#3}{l}{\small\sl continúa desde la página anterior}\\
      \toprule
      #5
      \otoprule
    }
    \tabletail{
      \hline
      \multicolumn{#3}{r}{\small\sl continúa en la página siguiente}\\
    }
    \tablelasttail{
      \hline
    }
    \bottomcaption{#1}
    \begin{xtabular}{#2}
      #6
      \bottomrule
    \end{xtabular}
    \label{tabla:#4}
  \end{center}
}

%
% Nuevo comando para tablas grandes sin cabecera.
\newcommand{\tablaSinCabecera}[5]{%
  \begin{center}
    \tablefirsthead{
      \toprule
    }
    \tablehead{
      \multicolumn{#3}{l}{\small\sl continúa desde la página anterior}\\
      \hline
    }
    \tabletail{
      \hline
      \multicolumn{#3}{r}{\small\sl continúa en la página siguiente}\\
    }
    \tablelasttail{
      \hline
    }
    \bottomcaption{#1}
  \begin{xtabular}{#2}
    #5
   \bottomrule
  \end{xtabular}
  \label{tabla:#4}
  \end{center}
}



\definecolor{cgoLight}{HTML}{EEEEEE}
\definecolor{cgoExtralight}{HTML}{FFFFFF}

%
% Nuevo comando para tablas grandes sin cabecera.
\newcommand{\tablaSinCabeceraConBandas}[5]{%
  \begin{center}
    \tablefirsthead{
      \toprule
    }
    \tablehead{
      \multicolumn{#3}{l}{\small\sl continúa desde la página anterior}\\
      \hline
    }
    \tabletail{
      \hline
      \multicolumn{#3}{r}{\small\sl continúa en la página siguiente}\\
    }
    \tablelasttail{
      \hline
    }
    \bottomcaption{#1}
    \rowcolors[]{1}{cgoExtralight}{cgoLight}

  \begin{xtabular}{#2}
    #5
   \bottomrule
  \end{xtabular}
  \label{tabla:#4}
  \end{center}
}




\graphicspath{ {./img/} }

% Capítulos
\chapterstyle{bianchi}
\newcommand{\capitulo}[2]{
	\setcounter{chapter}{#1}
	\setcounter{section}{0}
	\setcounter{figure}{0}
	\setcounter{table}{0}
	\chapter*{#2}
	\addcontentsline{toc}{chapter}{#2}
	\markboth{#2}{#2}
}

% Apéndices
\renewcommand{\appendixname}{Apéndice}
\renewcommand*\cftappendixname{\appendixname}

\newcommand{\apendice}[1]{
	%\renewcommand{\thechapter}{A}
	\chapter{#1}
}

\renewcommand*\cftappendixname{\appendixname\ }

% Formato de portada
\makeatletter
\usepackage{xcolor}
\newcommand{\tutor}[1]{\def\@tutor{#1}}
\newcommand{\course}[1]{\def\@course{#1}}
\definecolor{cpardoBox}{HTML}{E6E6FF}
\def\maketitle{
  \null
  \thispagestyle{empty}
  % Cabecera ----------------
\noindent\includegraphics[width=\textwidth]{cabecera}\vspace{1cm}%
  \vfill
  % Título proyecto y escudo informática ----------------
  \colorbox{cpardoBox}{%
    \begin{minipage}{.8\textwidth}
      \vspace{.5cm}\Large
      \begin{center}
      \textbf{TFG del Grado en Ingeniería Informática}\vspace{.6cm}\\
      \textbf{\LARGE\@title{}}
      \end{center}
      \vspace{.2cm}
    \end{minipage}

  }%
  \hfill\begin{minipage}{.20\textwidth}
    \includegraphics[width=\textwidth]{escudoInfor}
  \end{minipage}
  \vfill
  % Datos de alumno, curso y tutores ------------------
  \begin{center}%
  {%
    \noindent\LARGE
    Presentado por \@author{}\\ 
    en Universidad de Burgos --- \@date{}\\
    Tutor: \@tutor{}\\
  }%
  \end{center}%
  \null
  \cleardoublepage
  }
\makeatother


% Datos de portada
\title{título del TFG \\Documentación Técnica}
\author{nombre alumno}
\tutor{nombre tutor}
\date{\today}

\begin{document}

\maketitle



\cleardoublepage



%%%%%%%%%%%%%%%%%%%%%%%%%%%%%%%%%%%%%%%%%%%%%%%%%%%%%%%%%%%%%%%%%%%%%%%%%%%%%%%%%%%%%%%%



\frontmatter


\clearpage

% Indices
\tableofcontents

\clearpage

\listoffigures

\clearpage

\listoftables

\clearpage

\mainmatter

\appendix

\apendice{Plan de Proyecto Software}

\section{Introducción}

En este anexo se va a exponer la planificación temporal junto con la viabilidad legal y económica del proyecto.

\section{Planificación temporal}

\subsection{Introducción}
Para gestionar el proyecto se ha utilizado Github junto con Zube, primero disponemos de una parte temporal la cual no se disponía del repositorio Github en la cual se prueba la viabilidad del proyecto, después tenemos el desarrollo del proyecto junto con los sprints y por ultimo una parte que no se encuentra en Github correspondiente a la implementación de la documentación. En este flujo podemos observar como los comits han sido realizados.
\imagen{flujoCommits}{Flujo de commits}
\subsection{Sprints}
 \begin{itemize}
    \item Sprint previo 20/2/2024 – 22/3/2024: probar la viabilidad realizando un pequeño test que tenga más de 2 preguntas y varias preguntas cada una, además de recorrer el resultado y mostrarlo por pantalla, creación de usuarios de manera forzada con un array y creación de pestaña para logearse.
    \item Sprint 1 22/3/2024 – 3/4/2024: Implementación de la documentación. Añadir sesiones al test. Mostrar imágenes en el test. Mejorar la visualización de las preguntas.Se aplica la librería flask-session para poder operar con las sesiones y utilización de Bootstrap para mejorar la interfaz de la página web.
    \item Sprint 2 3/4/2024 – 17/4/2024: Base de datos con usuarios y test resueltos. Implementación de registro. En caso de no estar logeado hacer el test como un invitado. Guardar valores de los test con su nombre de usuario, fecha y puntuación. Opción de recuperación de contraseña. Ir avanzando en documentación. Utilización de flask-sqlalchemy y flask-migrate para el tratamiento de las BBDD.
    \item Sprint 3 17/4/2024 – 2/5/2024: Documentación (general y manual de instalación). Corregir fallo con el usuario invitado. Implementar campos restantes en el registro. Formulario para resetear la contraseña. Implementar parte inicial al cuestionario. Mostrar usuario que está conectado en la página principal. Creación del manual de programador, intento de realización del cambio de contraseña utilizando verificación de dos pasos, implementación de la pestaña previa si se realiza el test como invitado.
    \item Sprint 4 2/5/2024 – 15/5/2024: Creación de administradores que puedan ver las tablas resultados y usuarios. Implementación del otro tipo de test. Cambiar imágenes del test y añadir descripción. Mejora del cambio de contraseña. Avanzar en la documentación. Creación del rol administrador con permiso para ver los test y los usuarios de las BBDD, creación del cuestionario y adaptación de los valores a los originales aparte de rediseñar la ruta de las funciones para el invitado, intento de cambiar el cambio de contraseña.
    \item Sprint 5 15/5/2024 – 29/5/2024: Mejora del cambio de contraseña. Pulimiento de estética. Cambio de estética general y cambio del tratamiento del cambio de contraseña utilizando una pregunta secreta elegida aleatoriamente.
    \item Sprint 6 29/5/2024 – 6/6/2024: Memoria e implementación opción descargar resultados para administrador. Opción para descargar en un csv los resultados de los tests y avance en la documentación.
    \item Sprint Final 6/6/2024 – /7/2024: terminar memoria y anexos, comentar código, cambios mínimos y utilización de docker.
\end{itemize}
\section{Estudio de viabilidad}
En esta parte del apéndice se va a realizar un estudio sobre la viabilidad económica y legal del proyecto.
\subsection{Viabilidad económica}
A la hora de estudiar la posible viabilidad económica, se deben tener en cuenta varios aspectos como el gasto de los dispositivos utilizados, licencias e infraestructura.
\subsubsection{Costes generales}
En este apartado se incluirá lo relacionado con el uso del hardware y software del proyecto. Todo hardware tiene un coste no como algunas licencias de software que son open source, pero se tendrán en cuenta aquellas que no.
\begin{enumerate}
    \item \textbf{Costes del hardware:} He utilizado mi ordenador personal que tiene las siguientes especificaciones y en su momento tuvo un coste de 800€.
    \begin{itemize}
        \item CPU Intel(R) Core(TM) i7-7700HQ
        \item 8 GB de RAM
        \item GPU NVIDIA GeForce GTX 1050
    \end{itemize}
    La media de utilización de un portátil suele ser sobre los 4 años, por lo que el coste anual es:
    \begin{equation}
        \text{Coste anual} = \frac{800\,\text{€}}{4\,\text{años}} = 200\,\text{€/año}
    \end{equation}
    \item \textbf{Costes del software:} Para el proyecto se ha utilizado Docker el cual dispone de una licencia estándar gratuita, pero en caso de querer dar un uso mas avanzado es recomendable tener la versión pro como mínimo que tiene un coste de 5 \text{€} mensuales.
     \begin{equation}
        \text{Coste anual} = 5\,\text{€ x 12 meses} = 70\,\text{€/año}
    \end{equation}
    Por otra parte el uso de Github y Visual Studio Code son licencias gratuitas por lo que se tendrá en cuenta. Por último el sistema operativo utilizado que es Windows 10 ya no se encuentra en venta por lo que utilizaremos como referencia el Windows 11 Pro el cual su licencia cuesta 259€.
    \begin{equation}
        \text{Coste anual} = \frac{259\,\text{€}}{4\,\text{años}} = 64.75\,\text{€/año}
    \end{equation}
\end{enumerate}
Teniendo en cuenta todos los costes tendríamos el resultado de la siguiente tabla:
\\
\\
\begin{table}[H]
    \centering
    \begin{tabular}{@{}|l|r|@{}}
        \toprule
        Recurso & Precio/año \\ \midrule
        Ordenador & 200,00 € \\ 
        Licencia Docker & 70,00 € \\ 
        Windows 11 Pro & 64,75 € \\ \midrule
        Total: & 334,75 € \\ \bottomrule
    \end{tabular}
    \caption{Coste anual por recursos hardware y software.}
\end{table}

\subsubsection{Costes de infraestructura y personal}
Hay que tener en cuenta que a la hora de realizar un trabajo es necesario tener una sede, por lo cual se debe realizar un alquiler o compra de un local, el cual el precio medio de alquiler en España ronda a los 300€/mes.
\begin{equation}
    \text{Coste anual} = 300\,\text{€ x 12 meses} = 3600\,\text{€/año}
\end{equation}
En cuanto al personal el único trabajador soy yo, lo cual el sueldo correspondiente para un desarrollador Junior es de media 21.000€ anuales.
\begin{equation}
    \text{Salario bruto mensual} = \frac{21000\,\text{€}}{14\,\text{pagas}} = 1500\,\text{€}
\end{equation}
Por lo cual el costa anual sera:
\begin{equation}
    \text{Coste anual} = 1500\,\text{€ x 12 meses} = 18000\,\text{€/año}
\end{equation}
Teniendo en cuenta todos los costes tendríamos el resultado de la siguiente tabla:
\\
\\
\begin{table}[H]
    \centering
    \begin{tabular}{@{}|l|r|@{}}
         \toprule
        Recurso & Precio \\ \midrule
        Coste del software & 3600 € \\ 
        Coste de infraestructura y personal & 18000 € \\ \midrule
        Total: & 21600 € \\ \bottomrule 
    \end{tabular}
    \caption{Coste anual por personal e infraestructura.}
\end{table}

Una vez ya tenemos todos los costes a tener en cuenta, tenemos que buscar la manera en la cual sea sostenible y se puedan tener beneficios. Por lo que primero vamos a determinar un precio anual con el que se pueda contratar el servicio de RiskReal con las posibles modificaciones y mantenimiento.
Al ser una suscripción anual vamos a poner un precio estándar de 1000 € por lo que vamos a calcular el total de clientes necesarios que se deberán tener para pasar el margen de beneficio:
\\
\\
\begin{table}[H]
    \centering
    \begin{tabular}{@{}|l|r|@{}}
         \toprule
        Recurso & Precio/año \\ \midrule
        Infraestructura & 334,75 € \\
        Personal & 21600,00 € \\ \midrule
        Total: & 21934,75 € \\ \bottomrule 
    \end{tabular}
    \caption{Coste anual total}
\end{table}

\begin{equation}
    \text{Numero clientes} = \frac{21934.75\,\text{€}}{1000\,\text{€}} = 21.93\,\text{ clientes}
\end{equation}

Realizado el estudio tenemos que para que el proyecto sea viable debemos tener como mínimo un total de 22 clientes anuales, para tener una economía sostenible con alguna ganancia. 
\subsection{Viabilidad legal}
Para evaluar la viabilidad legal de un proyecto de software, es crucial considerar varios factores como las licencias de software utilizadas, y las leyes y regulaciones aplicables. A continuación, se presenta un análisis detallado de la viabilidad legal del proyecto.

\subsubsection{Licencias}
\begin{table}[]
    \centering
    \begin{tabular}{@{}|l|l|l|@{}}
        \toprule
        \textbf{Librería} & \textbf{Versión} & \textbf{Licencia} \\ \hline
        alembic & 1.13.1 & MIT License \\ \midrule
        blinker &1.8.2 & MIT License \\ \midrule
        click & 8.1.7 & BSD-3-Clause License \\ \midrule
        colorama & 0.4.6 & BSD-3-Clause License \\ \midrule
        Flask & 3.0.3 & BSD-3-Clause License \\ \midrule
        Flask-Migrate & 4.0.7 & MIT License \\ \midrule
        Flask-SQLAlchemy & 3.1.1 & BSD-3-Clause License \\ \midrule
        greenlet & 3.0.3 & MIT License \\ \midrule
        install & 1.3.5 & MIT License \\ \midrule
        itsdangerous & 2.2.0 & BSD-3-Clause License \\ \midrule
        Jinja2 & 3.1.4 & BSD-3-Clause License \\ \midrule
        Mako & 1.3.5 & MIT License \\ \midrule
        MarkupSafe & 2.1.5 & BSD-3-Clause License \\ \midrule
        pip & 24.0 & MIT License \\ \midrule
        python-dotenv & 1.0.1 & BSD-3-Clause License \\ \midrule
        setuptools & 65.5.0 & MIT License \\ \midrule
        SQLAlchemy & 2.0.9 & MIT License \\ \midrule
        typing\_extensions &4.12.0 & Python Software Foundation License \\ \midrule
        Werkzeug & 3.0.3 & BSD-3-Clause License \\ \bottomrule 
    \end{tabular}
    \caption{LIbrerias del proyecto con su version y licencia}
\end{table}

\begin{itemize}
    \item \textbf{MIT License} ~\cite{wiki:MIT}: Permite a los usuarios hacer casi cualquier cosa con el proyecto, incluyendo usarlo, copiarlo, modificarlo, fusionarlo, publicarlo, distribuirlo, sublicenciarlo y vender copias del software, siempre que incluyan el aviso de derechos de autor y la exención de responsabilidad en todas las copias o partes sustanciales del software.
    \item \textbf{BSD-3-Clause License} ~\cite{wiki:BSD}: Similar a la MIT, pero con una cláusula adicional que impide el uso del nombre de los autores para promocionar productos derivados sin permiso.
    \item \textbf{Python Software Foundation License} ~\cite{wiki:PSFL}: Es una licencia específica para Python que es similar en espíritu a las licencias MIT y BSD. Es permisiva y permite el uso, copia, modificación y distribución del software.
\end{itemize}
\subsubsection{Conclusión}
El proyecto es legalmente viable teniendo en cuenta que las licencias utilizadas permiten la distribución siempre que se tengan en cuenta la propiedad intelectual del código. No se encuentran barreras legales para la continuación del desarrollo del proyecto, al utilizar solo licencias permisivas.

Cabe destacar que es aconsejable la consulta con un abogado especializado para estar seguro de seguir correctamente toda la jurisdicción  correspondiente.

Como resultado la licencia que utilizaria mi proyecto seria de tipo \textbf{BSD-3-Clause License}.
\apendice{Especificación de Requisitos}

\section{Introducción}

Una muestra de cómo podría ser una tabla de casos de uso:

% Caso de Uso 1 -> Consultar Experimentos.
\begin{table}[p]
	\centering
	\begin{tabularx}{\linewidth}{ p{0.21\columnwidth} p{0.71\columnwidth} }
		\toprule
		\textbf{CU-1}    & \textbf{Ejemplo de caso de uso}\\
		\toprule
		\textbf{Versión}              & 1.0    \\
		\textbf{Autor}                & Alumno \\
		\textbf{Requisitos asociados} & RF-xx, RF-xx \\
		\textbf{Descripción}          & La descripción del CU \\
		\textbf{Precondición}         & Precondiciones (podría haber más de una) \\
		\textbf{Acciones}             &
		\begin{enumerate}
			\def\labelenumi{\arabic{enumi}.}
			\tightlist
			\item Pasos del CU
			\item Pasos del CU (añadir tantos como sean necesarios)
		\end{enumerate}\\
		\textbf{Postcondición}        & Postcondiciones (podría haber más de una) \\
		\textbf{Excepciones}          & Excepciones \\
		\textbf{Importancia}          & Alta o Media o Baja... \\
		\bottomrule
	\end{tabularx}
	\caption{CU-1 Nombre del caso de uso.}
\end{table}

\section{Objetivos generales}

\section{Catálogo de requisitos}

\section{Especificación de requisitos}



\apendice{Especificación de diseño}

\section{Introducción}

\section{Diseño de datos}

\section{Diseño procedimental}

\section{Diseño arquitectónico}



\apendice{Documentación técnica de programación}

\section{Introducción}
En este apartado se implementaran  los conocimientos necesarios para la instalación y funcionamiento  de la aplicación. Esto incluye aspectos como la diferente estructura de los directorios, instalación y uso del entorno de desarrollo, y la propia compilación, instalación y ejecución del proyecto.
\section{Estructura de directorios}
Todo el código se encuentra en \textbf{main.py}, el resto de ficheros de la raíz son ficheros de las páginas y miscelánea utilizada, se modificaran aquellos que sean necesario respecto a su uso en el código.
\begin{enumerate}
    \item \textbf{json/}: Contiene todos los ficheros json con los datos de los test.    
    \item \textbf{static/assets/}: Contiene imágenes e iconos utilizados en los html.
    \item \textbf{static/css/}: Contiene los diferentes archivos css para modificar la vista de los archivos html.
    \item \textbf{static/gifs/}: Contiene los diferentes archivos gifs que se introducen en los propios html.
    \item \textbf{static/img/}: Contiene las carpetas de los diferentes test con las imágenes utilizadas en los json.
    \item \textbf{static/js/}: Contiene ficheros de java script utilizados por bootstrap.
    \item \textbf{templates/}: Contiene los archivos html de las diferentes páginas de la aplicación.
\end{enumerate} 
\section{Manual del programador}
En este apartado tendremos los pasos para instalar los entornos sobre el que se trabaja.
\subsection{Instalación de Visual Studio Code}
En caso de no tener \textbf{Visual Studio Code} en el sistema, nos descargaremos la versión correspondiente para nuestro sistema operativo \href{https://code.visualstudio.com/}{página oficial} una vez descargado el ejecutable lo lanzamos y completamos la instalación.
\subsection{Instalación de Python}
Para poder tener python en nuestro ordenador deberemos descargar la versión que se encuentra en la \href{https://www.python.org/downloads/}{página oficial} correspondiente para nuestro sistema operativo.
Una vez lo tengamos ya descargado inicializamos él .exe y marcamos la opción que dice textualmente  "Add python.exe to PATH", realizamos la instalación y una vez se termine seguiremos con el siguiente paso.
En caso de no tener \textbf{Python} en \textbf{Visual Studio Code}, seguiremos los siguientes pasos:
\begin{enumerate}
    \item En la barra de la izquierda de \textbf{Visual Studio Code} seleccionamos la opción de Extensions.
    \item En el buscador introduccimos \textbf{Python}
    \item Seleccionamos la primera opción la cual tiene como desarrollador Microsoft y la instalamos.
\end{enumerate}
\subsection{Instalación de Git}
Para poder llevar la versión del proyecto al día utilizaremos git el cual se vinculara solo con \textbf{Visual Studio Code}, para instalarlo nos dirigiremos a la \href{https://git-scm.com/download/win}{página oficial} y nos descargamos el instalador, una vez ya lo tengamos ejecutamos él .exe y lo instalamos en nuestro dispositivo.
\subsection{Instalación de Docker}
\label{InstDocker}
Para poder descargar correctamente \textbf{Docker} hay que seguir correctamente estos \href{https://docs.docker.com/desktop/install/windows-install/}{pasos}, como referencia se pueden seguir los pasos del siguiente \href{https://www.youtube.com/watch?v=ZO4KWQfUBBc}{video} el cual realiza los pasos de forma secuencial y así podemos comprobar que realizamos correctamente la instalación.
\subsubsection{Descarga del proyecto}
Para poder tener los ficheros fuentes, en el lugar que se desea hacemos click derecho y selecionamos \textbf{Open Git Bash here}, una vez se nos abra la terminal introduciremos el siguiente comando para descargar los ficheros fuentes directamente del repositorio:

\begin{itemize}
    \item git clone https://github.com/muf1002/TFG-RiskReal.git
\end{itemize}

\imagen{clone}{Comando para clonar el repositorio}

Posteriormente desde \textbf{Visual Studio Code} damos la opción de open folder o nos dirigimos en la barra de tareas superior donde se encuentra también, seleccionamos la carpeta fuente tras clonarlo y de esta manera tendríamos los ficheros en el entorno.

\section{Compilación, instalación y ejecución del proyecto}
Una vez tengamos los pasos previos seguiremos las siguientes indicaciones para completar la instalación y poder ejecutarlo, es recomendable realizar la instalación teniendo los permisos de administrador.
\begin{enumerate}
    \item Creamos un fichero \textbf{.env} al nivel de la raíz del proyecto para poder declarar variables del entorno, en el cual introduciremos el valor sobre CLAVE\_ADMIN=\textbf{clave} siendo clave el valor deseado para este.
    \imagen{env}{Creación del fichero .env}
    \item Abrimos una terminal para poder ejecutar los comandos necesarios 
    \item py -3 -m venv RiskReal: creamos un entorno virtual siendo \textbf{RiskReal} el identificador para saber que estamos trabajando sobre \textbf{RiskReal}.
    \imagen{venv}{Creación del entorno virtual}
    \item RisReal\textbackslash Scripts\textbackslash activate: una vez creado el entorno lo activamos.
    \imagen{activate}{Activación del entorno virtual}
    \item pip install -r requirements.txt: una vez activado el entorno instalaremos todos los elementos necesarios.
    \imagen{installRequirements}{Instalación de las librerias}
    \item \$env:FLASK\_APP = "main": con este comando estableceremos como fichero fuente para la aplicación \textbf{FLASK} el fichero \textbf{main.py}, en caso de dar algun tipo de error hay que ejecutar la consola como administrador.
    \imagen{entorno}{Vinculación del fichero main.py con Flask}
    \item flask db init: una vez ya tenemos todo configurado creamos los ficheros principales para \textbf{flask\_migrate}.
    \imagen{init}{Inicialización de migrate}
    \item flask db migrate: una vez creados se generan las instancias de las tablas.
    \imagen{migrate}{Instancia de las tablas}
    \item flask db upgrade: introducion este comando generaremos las tablas y en caso de que hayan sido modificadas se ejecutara para aplicar los cambios de estas antes de ejecutar el proyecto.
    \imagen{upgrade}{Generación de las tablas}
    \item flask \text{-}-app main run \text{-}-debug: por último para ejecutar el propio proyecto se lanzara el comando con la opción de \textbf{debug} para tener una traza de la ejecución del proyecto.
    \imagen{mainRun}{Lanzamiento de la aplicación}
\end{enumerate}

Una vez que entremos a \textbf{Visual Studio Code} y ya tengamos la instalación previa de otro momento, solo tenemos que activar el entorno (paso 4) y una vez ya nos encontremos en el entorno ejecutaremos el comando del paso 6, así si queremos ejecutar la aplicación de nuevo solo tenemos que realizar el paso 10 o si en su caso se modifica algún campo de las tablas se realizaran los pasos 8 y 9 para poner en contexto la aplicación.

Para poder crear un contenedor con Docker seguiremos los siguientes pasos:
\begin{enumerate}
    \item Seguimos todos los pasos anteriores hasta el paso numero 10 ya que no queremos ejecutar la aplicación, sino tener los archivos necesarios para que RiskReal pueda ser lanzado.
    \item Para poder realizar cualquier acción con Docker debemos tener Docker Desktop en funcionamiento, por lo que si no lo tenemos ejecutando lo abrimos.
    \item Una vez ya tenemos Docker Desktop en funcionamiento ejecutamos el siguiente comando para crear una imagen: \textbf{docker build -t riskreal .}, tardara unos segundos en ejecutarse.
    \imagen{dockerBuild}{Creación de la imagen de docker}
    \item Ahora para crear el contenedor el cual se ejecutara realizamos el siguiente comando \textbf{docker run -p 5000:5000 \text{-}-env-file .env riskreal}
    \imagen{dockerRun}{Creación y lanzamiento del contenedor}
    \item Una vez ejecutado podemos acceder a la página web ya sea con la primera dirección que se muestra o con la dirección (localhost:5000).
\end{enumerate}

Una vez ya tengamos el contenedor si nos dirigimos al Docker Desktop podemos ver el contenedor creado, para ejecutarlo solo tendríamos que darle al botón de play que aparece en la parte de la derecha, y posteriormente clicamos en el puerto el cual nos abrirá la página web. No es necesario tener abierto Visual Studio Code para poder lanzar correctamente la aplicación.
\imagen{contenedor}{Contenedor generado en Docker Desktop}
\section{Pruebas del sistema}
En cuanto a las pruebas a realizar para asegurar el correcto funcionamiento del \textbf{proyecto} a medida que se implementaba un nuevo requisito o se modificaba, al momento de comprobar que se realizase la nueva funcionalidad se utiliza el modo debug para poder recoger los fallos y asi poder correguirlos.

A la hora de realizar una prueba conjunta se ha realizado de la misma manera siguiendo los pasos que se deberían seguir para comprobar el funcionamiento completo.
 
\apendice{Documentación de usuario}

\section{Introducción}

\section{Requisitos de usuarios}

\section{Instalación}

\section{Manual del usuario}



\apendice{Anexo de sostenibilización curricular}

\section{Introducción}
Este anexo incluirá una reflexión personal del alumnado sobre los aspectos de la sostenibilidad que se abordan en el trabajo.
Se pueden incluir tantas subsecciones como sean necesarias con la intención de explicar las competencias de sostenibilidad adquiridas durante el alumnado y aplicadas al Trabajo de Fin de Grado.

Más información en el documento de la CRUE \url{https://www.crue.org/wp-content/uploads/2020/02/Directrices_Sosteniblidad_Crue2012.pdf}.

Este anexo tendrá una extensión comprendida entre 600 y 800 palabras.



\bibliographystyle{plain}
\bibliography{bibliografiaAnexos}

\end{document}
